\section{Математическое описание}

\subsection{Определение}

\texttt{bf16} --- формат хранения числа с точкой в 16~битах: один бит
знака $S$, 8 бит смещённой экспоненты $E$ (смещение $2^{k-1} - 1 = 2^{8-1}-1 = 127$) и семь битов
мантиссы $M$. Число хранится в нормализованном виде
$1.M \cdot 2^{E-127}$, при этом единица перед точкой явно не хранится.
Вещественное число вычисляется по формуле:
\[
    x = (-1)^s \cdot (1 + M) \cdot 2^{E-127}, \,\, \text{где}\,\, M = 0.xxxx
\]

\subsection{Особые значения}

Основная формула применяется только при $1 \leq E \leq 254$. Значения
экспоненты, состоящие только из нулей или единиц, имеют специальный смысл:
\[
\begin{array}{c|c|c}
E & M & \text{значение} \\
\hline
00000000 & 0000000 & \text{ноль} \\
00000000 & M \neq 0000000 & \text{денормализованное число} \\
11111111 & 0000000 & \text{бесконечность} \\
11111111 & M \neq 0000000 & \text{\texttt{NaN}}
\end{array}
\]

Ноль имеет вид $S \;00000000\;0000000$. Бит знака позволяет различать
$+0$ и $-0$, однако их численное значение равно нулю:
\[
    x = (-1)^S \cdot 0.
\]

Если $E=0$, а мантисса не равна нулю, получается денормализованное очень
маленькое число. У него нет неявной единицы перед мантиссой:
\[
    x = (-1)^S \cdot M \cdot 2^{1-127}.
\]

Бесконечность имеет вид $S\;11111111\;0000000$. Значение бита знака
определяет, будет это $+\infty$ или $-\infty$. Такая запись получается,
когда число не помещается в диапазон формата.

Значение \texttt{NaN} имеет вид $S\;11111111\;M$, где хотя бы один бит
мантиссы равен единице. \texttt{NaN} означает неопределённый результат,
например $0/0$, и не является обычным числом.

\subsection{Переводы}

\subsubsection{Перевод из десятичной системы счисления в bf16}
Сначала переводим целую часть числа в двоичничную делением на два, дробную часть умножением на два(на каждом шаге умножения целая часть дробного числа очередной бит, остатов идет на следующий шаг). То, что получили сдвигаем до момента $1.xxxx * 2^x$, после считаем экспоненту $E= 127 + x$, посе чего 7 битов мантиссы округляем по следующему за ним биту(округление к ближайшему).

\subsubsection{Перевод из bf16 в десятичную систему счисления}
Разбиваем 16 бит на знак $S$, экспоненту $E$, мантиссу $M$, после чего значение считаем по формуле выше(с учетом крайних случаев)
